% Options for packages loaded elsewhere
\PassOptionsToPackage{unicode}{hyperref}
\PassOptionsToPackage{hyphens}{url}
\documentclass[
  ignorenonframetext,
]{beamer}
\newif\ifbibliography
\usepackage{pgfpages}
\setbeamertemplate{caption}[numbered]
\setbeamertemplate{caption label separator}{: }
\setbeamercolor{caption name}{fg=normal text.fg}
\beamertemplatenavigationsymbolsempty
% remove section numbering
\setbeamertemplate{part page}{
  \centering
  \begin{beamercolorbox}[sep=16pt,center]{part title}
    \usebeamerfont{part title}\insertpart\par
  \end{beamercolorbox}
}
\setbeamertemplate{section page}{
  \centering
  \begin{beamercolorbox}[sep=12pt,center]{section title}
    \usebeamerfont{section title}\insertsection\par
  \end{beamercolorbox}
}
\setbeamertemplate{subsection page}{
  \centering
  \begin{beamercolorbox}[sep=8pt,center]{subsection title}
    \usebeamerfont{subsection title}\insertsubsection\par
  \end{beamercolorbox}
}
% Prevent slide breaks in the middle of a paragraph
\widowpenalties 1 10000
\raggedbottom
\AtBeginPart{
  \frame{\partpage}
}
\AtBeginSection{
  \ifbibliography
  \else
    \frame{\sectionpage}
  \fi
}
\AtBeginSubsection{
  \frame{\subsectionpage}
}
\usepackage{iftex}
\ifPDFTeX
  \usepackage[T1]{fontenc}
  \usepackage[utf8]{inputenc}
  \usepackage{textcomp} % provide euro and other symbols
\else % if luatex or xetex
  \usepackage{unicode-math} % this also loads fontspec
  \defaultfontfeatures{Scale=MatchLowercase}
  \defaultfontfeatures[\rmfamily]{Ligatures=TeX,Scale=1}
\fi
\usepackage{lmodern}
\ifPDFTeX\else
  % xetex/luatex font selection
\fi
% Use upquote if available, for straight quotes in verbatim environments
\IfFileExists{upquote.sty}{\usepackage{upquote}}{}
\IfFileExists{microtype.sty}{% use microtype if available
  \usepackage[]{microtype}
  \UseMicrotypeSet[protrusion]{basicmath} % disable protrusion for tt fonts
}{}
\makeatletter
\@ifundefined{KOMAClassName}{% if non-KOMA class
  \IfFileExists{parskip.sty}{%
    \usepackage{parskip}
  }{% else
    \setlength{\parindent}{0pt}
    \setlength{\parskip}{6pt plus 2pt minus 1pt}}
}{% if KOMA class
  \KOMAoptions{parskip=half}}
\makeatother
\usepackage{longtable,booktabs,array}
\newcounter{none} % for unnumbered tables
\usepackage{calc} % for calculating minipage widths
\usepackage{caption}
% Make caption package work with longtable
\makeatletter
\def\fnum@table{\tablename~\thetable}
\makeatother
\setlength{\emergencystretch}{3em} % prevent overfull lines
\providecommand{\tightlist}{%
  \setlength{\itemsep}{0pt}\setlength{\parskip}{0pt}}
% Auto-generated by infrastructure/rendering/_pdf_latex_helpers.py::extract_math_font_preamble.
% Minimal header passed to Pandoc via -H so Beamer slides pick up the same math font
% as the combined-PDF path. Adjust the manuscript's preamble.md to change the font globally.
\usepackage{unicode-math}
\setmathfont{latinmodern-math.otf}

% Slide layout defaults for warning-clean scientific decks.
\usepackage{etoolbox}
\IfFileExists{xurl.sty}{\usepackage{xurl}}{}
\IfFileExists{seqsplit.sty}{\usepackage{seqsplit}}{\newcommand{\seqsplit}[1]{#1}}
\protected\def\breakseq#1{\seqsplit{#1}}
\protected\def\breaktt#1{\begingroup\ttfamily\seqsplit{#1}\endgroup}
\setlength{\emergencystretch}{6em}
\tolerance=5000
\hbadness=10000
\hfuzz=1pt
\setlength{\tabcolsep}{2pt}
\AtBeginEnvironment{longtable}{\tiny\renewcommand{\arraystretch}{0.86}\setlength{\tabcolsep}{1pt}}
\AtBeginEnvironment{tabular}{\tiny\renewcommand{\arraystretch}{0.86}\setlength{\tabcolsep}{1pt}}

% Natbib and cross-reference fallbacks — slides don't load natbib
% or cleveref, but combined-PDF manuscript prose may emit these
% commands. The fallback renders citations as a bracketed key list
% and cross-references as detokenized labels so slides don't fail on
% undefined control sequences or raw underscores. \providecommand is
% a no-op if packages load later.
\providecommand{\citep}[1]{[#1]}
\providecommand{\citet}[1]{#1}
\providecommand{\citealp}[1]{#1}
\providecommand{\citeauthor}[1]{#1}
\providecommand{\citeyear}[1]{#1}
\providecommand{\cref}[1]{\texttt{\detokenize{#1}}}
\providecommand{\Cref}[1]{\texttt{\detokenize{#1}}}
\usepackage{bookmark}
\IfFileExists{xurl.sty}{\usepackage{xurl}}{} % add URL line breaks if available
\urlstyle{same}
% fallback for those not using the hyperref driver hyperxmp:
\makeatletter
\@ifundefined{xmpquote}{\newcommand{\xmpquote}[1]{#1}}{}
\makeatother
\hypersetup{
  hidelinks,
  pdfcreator={LaTeX via pandoc}}

\author{\texorpdfstring{}{}}
\date{}

\begin{document}

\section{Introduction: Lexicon as Data and Manuscript as Build
Artifact}\label{introduction-lexicon-as-data-and-manuscript-as-build-artifact}

\begin{frame}[fragile,allowframebreaks]{Introduction: Lexicon as Data
and Manuscript as Build Artifact}
Mad Lib style generation is usually treated as a toy because it
foregrounds the visible blank rather than the source of the replacement.
In a research pipeline, the blank is not the hard part. The hard part is
making every replacement reviewable, deterministic, and honest about
what it can support. \texttt{template\_madlib} turns that constraint
into the subject of the exemplar.

The project treats nouns such as pipeline, protocol, section, lexicon
and verbs such as hydrate, condition, bind, bind as versioned data.
Changing a lexicon entry is therefore closer to changing an input table
than editing prose in place. That distinction matters because the
generated manuscript can be rerendered, diffed, and validated without
asking the reader to trust an invisible drafting session.

The introduction is configured to separate playful Mad Lib syntax from
research claims, identify drift between prose and source data, frame
configuration as an inspectable dataset, and position conditional prose
as a reproducibility problem. Those moves keep the manuscript from
pretending to be an open-ended language model. It is a bounded template:
authors declare categories, slots, section switches, method steps, and
claim boundaries; source code transforms those declarations into
manuscript bodies and evidence tables.

This exemplar is useful for protocols, educational scaffolds, review
forms, templated reports, and other documents where conditional text is
unavoidable but should never become untraceable. The same pattern can be
extended with domain-specific validators while leaving the shared
rendering infrastructure untouched.
\end{frame}

\begin{frame}[allowframebreaks]{Contribution Ledger}
\protect\phantomsection\label{contribution-ledger}
{\def\LTcaptype{none} % do not increment counter
\begin{longtable}[]{@{}
  >{\raggedright\arraybackslash}p{(\linewidth - 2\tabcolsep) * \real{0.5000}}
  >{\raggedright\arraybackslash}p{(\linewidth - 2\tabcolsep) * \real{0.5000}}@{}}
\toprule\noalign{}
\begin{minipage}[b]{\linewidth}\raggedright
Claim
\end{minipage} & \begin{minipage}[b]{\linewidth}\raggedright
Boundary
\end{minipage} \\
\midrule\noalign{}
\endhead
A Mad Lib manuscript can remain reproducible when the lexicon is treated
as data. & Local exemplar claim; no live DOI or standalone release
implied. \\
Conditional IMRAD section bodies can be rendered without shared renderer
changes. & Local exemplar claim; no live DOI or standalone release
implied. \\
Large-grain manuscript variables can preserve author-readable source
files while still producing a complete manuscript. & Local exemplar
claim; no live DOI or standalone release implied. \\
Token provenance can connect playful lexical substitution to serious
publication hygiene. & Local exemplar claim; no live DOI or standalone
release implied. \\
A generated Methods section can be methodologically useful when protocol
rows, phases, figures, and validation gates share one config-owned
source. & Local exemplar claim; no live DOI or standalone release
implied. \\
Configured-field origin tracking can make loader defaults visible enough
for reviewers and downstream forks. & Local exemplar claim; no live DOI
or standalone release implied. \\
Pipeline-owned output regeneration can keep PDF, HTML, slides, data,
reports, and copied deliverables aligned without hand-editing generated
artifacts. & Local exemplar claim; no live DOI or standalone release
implied. \\
A generated-method exemplar can make review handoff auditable when the
review packet includes source config, data artifacts, figures,
validation results, and copy statistics. & Local exemplar claim; no live
DOI or standalone release implied. \\
Fork migration guidance can reduce overclaiming when it names the
source, test, validator, and evidence surfaces that must change before
domain use. & Local exemplar claim; no live DOI or standalone release
implied. \\
\bottomrule\noalign{}
\end{longtable}
}
\end{frame}

\end{document}
